\hypertarget{appendix-73425b25118e}{}
\section*{Shift operations}
\section{Shift operations}\label{shift-operations}

\textbf{Curriculum depth:} Developed · \textbf{First seen:} C001 · \textbf{Latest development:} C002

\subsection{Definition and mechanism}\label{definition-and-mechanism}

Relocate bit significance; unsigned shifts can correspond to multiply or floor-divide by powers of two within their validity range.

Logical right shift fills high positions with zero; arithmetic right shift replicates the sign bit. In C++23, signed right shift rounds toward negative infinity; signed integer division rounds toward zero. Thus −5\textgreater\textgreater1 gives −3 while −5/2 gives −2. A built-in shift count must be nonnegative and smaller than the promoted left operand's width.

\subsection{Prerequisites}\label{prerequisites}

\begin{itemize}
\tightlist
\item
  \hyperref[appendix-280445746493]{Bit significance} --- LSB means least significant bit, normally position 0 with weight 1.
\item
  \hyperref[appendix-d67a19507dc5]{Machine word} --- A processor's natural integer processing width is often called its word width, but the term is architecture dependent.
\end{itemize}

\subsection{Used by}\label{used-by}

\begin{itemize}
\item
  \hyperref[appendix-ac697f5c66e8]{Bit field} --- uses shift operations as part of its representation or reasoning.
\item
  \hyperref[appendix-a65290d66ed9]{Shift-and-add multiplication} --- uses shift operations as part of its representation or reasoning.
\item
  \hyperref[appendix-89985e6a2685]{Binary long division} --- uses shift operations as part of its representation or reasoning.
\item
  \hyperref[appendix-367b6abf460e]{Little endian} --- depends on this concept; see its prerequisite explanation.
\item
  \hyperref[appendix-0c67cd5bf30a]{Pages and page offsets} --- depends on this concept; see its prerequisite explanation.
\end{itemize}

\subsection{Recall and next depth}\label{recall-and-next-depth}

Explain shift operations without the label, work the example above, then state the assumption whose removal would change the result.

\textbf{Pending:} signed-shift rules, compiler lowering, rotate instructions, and overshift behavior.

\subsection{Source and navigation}\label{source-and-navigation}

\href{Continuous\%20Notes/Computing\%20Foundations\%20-\%20From\%20Bits\%20to\%20C++\%20and\%20Systems\%20APIs\#From\%20a\%20signal\%20to\%20an\%20interpretation}{C001 teaching chapter} · \href{Continuous\%20Notes/Computing\%20Foundations\%20-\%20From\%20Bits\%20to\%20C++\%20and\%20Systems\%20APIs\#Arithmetic\%20becomes\%20a\%20policy}{C002 development} · \hyperref[appendix-fd090d982a15]{Part definitions and mechanisms} · \href{}{Dependency Map (02-Foundations-II, appendix-4ad47bdd269a)}
